% !TEX TS-program = lualatex
% !TEX encoding = UTF-8

\documentclass[11pt, twoside, french]{book}

\input{common_headers}

\begin{document}

\feast{F5}{Feria V in Cœna Domini}{Feria V in Cœna Domini}{Feria V in Cœna Domini}{1}
\addcontentsline{toc}{part}{Feria V in Cœna Domini}


\feast{F5L1}{Lectio I}{Feria V in Cœna Domini}{Lectio I}{3}
\addcontentsline{toc}{chapter}{Lectio I}

\smalltitle{Tonus Romanus}
\addcontentsline{toc}{section}{Tonus Romanus}
\gscore{F5L1_roma}{}{}{}
\smalltitle{Tonus Silensis}
\addcontentsline{toc}{section}{Tonus Silensis}
\gscore{F5L1_silos}{}{}{}
\smalltitle{Tonus Toletanus}
\addcontentsline{toc}{section}{Tonus Toletanus}
\gscore{F5L1_toledo}{}{}{}

\feast{F5L2}{Lectio II}{Feria V in Cœna Domini}{Lectio II}{3}
\addcontentsline{toc}{chapter}{Lectio II}

\feast{F5L3}{Lectio III}{Feria V in Cœna Domini}{Lectio III}{3}
\addcontentsline{toc}{chapter}{Lectio III}

\feast{F5L4}{Lectio IV}{Feria V in Cœna Domini}{Lectio IV}{3}
\addcontentsline{toc}{chapter}{Lectio IV}

\feast{F5L5}{Lectio V}{Feria V in Cœna Domini}{Lectio V}{3}
\addcontentsline{toc}{chapter}{Lectio V}

\feast{F5L6}{Lectio VI}{Feria V in Cœna Domini}{Lectio VI}{3}
\addcontentsline{toc}{chapter}{Lectio VI}

\feast{F5L7}{Lectio VII}{Feria V in Cœna Domini}{Lectio VII}{3}
\addcontentsline{toc}{chapter}{Lectio VII}

\feast{F5L8}{Lectio VIII}{Feria V in Cœna Domini}{Lectio VIII}{3}
\addcontentsline{toc}{chapter}{Lectio VIII}

\feast{F5L9}{Lectio IX}{Feria V in Cœna Domini}{Lectio IX}{3}
\addcontentsline{toc}{chapter}{Lectio IX}




\feast{F6}{Feria VI in Parasceve}{Feria VI in Parasceve}{Feria VI in Parasceve}{1}
\addcontentsline{toc}{part}{Feria VI in Parasceve}

\feast{F6L1}{Lectio I}{Feria VI in Parasceve}{Lectio I}{3}
\addcontentsline{toc}{chapter}{Lectio I}

\feast{F6L2}{Lectio II}{Feria VI in Parasceve}{Lectio II}{3}
\addcontentsline{toc}{chapter}{Lectio II}

\feast{F6L3}{Lectio III}{Feria VI in Parasceve}{Lectio III}{3}
\addcontentsline{toc}{chapter}{Lectio III}

\feast{F6L4}{Lectio IV}{Feria VI in Parasceve}{Lectio IV}{3}
\addcontentsline{toc}{chapter}{Lectio IV}

\feast{F6L5}{Lectio V}{Feria VI in Parasceve}{Lectio V}{3}
\addcontentsline{toc}{chapter}{Lectio V}

\feast{F6L6}{Lectio VI}{Feria VI in Parasceve}{Lectio VI}{3}
\addcontentsline{toc}{chapter}{Lectio VI}

\feast{F6L7}{Lectio VII}{Feria VI in Parasceve}{Lectio VII}{3}
\addcontentsline{toc}{chapter}{Lectio VII}

\feast{F6L8}{Lectio VIII}{Feria VI in Parasceve}{Lectio VIII}{3}
\addcontentsline{toc}{chapter}{Lectio VIII}

\feast{F6L9}{Lectio IX}{Feria VI in Parasceve}{Lectio IX}{3}
\addcontentsline{toc}{chapter}{Lectio IX}




\feast{F7}{Sabbato Sancto}{Sabbato Sancto}{Sabbato Sancto}{1}
\addcontentsline{toc}{part}{Sabbato Sancto}

\feast{F7L1}{Lectio I}{Sabbato Sancto}{Lectio I}{3}
\addcontentsline{toc}{chapter}{Lectio I}

\feast{F7L2}{Lectio II}{Sabbato Sancto}{Lectio II}{3}
\addcontentsline{toc}{chapter}{Lectio II}

\feast{F7L3}{Lectio III}{Sabbato Sancto}{Lectio III}{3}
\addcontentsline{toc}{chapter}{Lectio III}

\feast{F7L4}{Lectio IV}{Sabbato Sancto}{Lectio IV}{3}
\addcontentsline{toc}{chapter}{Lectio IV}

\feast{F7L5}{Lectio V}{Sabbato Sancto}{Lectio V}{3}
\addcontentsline{toc}{chapter}{Lectio V}

\feast{F7L6}{Lectio VI}{Sabbato Sancto}{Lectio VI}{3}
\addcontentsline{toc}{chapter}{Lectio VI}

\feast{F7L7}{Lectio VII}{Sabbato Sancto}{Lectio VII}{3}
\addcontentsline{toc}{chapter}{Lectio VII}

\feast{F7L8}{Lectio VIII}{Sabbato Sancto}{Lectio VIII}{3}
\addcontentsline{toc}{chapter}{Lectio VIII}

\feast{F7L9}{Lectio IX}{Sabbato Sancto}{Lectio IX}{3}
\addcontentsline{toc}{chapter}{Lectio IX}


\newpage
\pagestyle{empty}
\renewcommand\contentsname{Index Generalis}
\tableofcontents
\thispagestyle{empty}

\end{document}

